\documentclass[12pt]{article}

\input{../../_assets/preambulo.tex}


\begin{document}

    % 1. Foto de fondo
    % 2. Título
    % 3. Encabezado Izquierdo
    % 4. Color de fondo
    % 5. Coord x del titulo
    % 6. Coord y del titulo
    % 7. Fecha

    
    \input{../../_assets/portada}
    \portadaExamen{ffccA4.jpg}{EDIP\\Examen IV}{EDIP. Examen IV}{MidnightBlue}{-8}{28}{2025-26}{Transcrito por Gemini}

    \begin{description}
        \item[Asignatura] Estadística Descriptiva e Introducción a la Probabilidad.
        \item[Curso Académico] 2025/2026.
        \item[Grado] Doble Grado en Ingeniería Informática y Matemáticas.
        \item[Grupo] Único
        \item[Profesor] Juan Antonio Maldonado Jurado.
        \item[Descripción] Convocatoria Ordinaria.
        \item[Fecha] 15 de junio de 2026.
        % \item[Duración] ---.
    
    \end{description}
    \newpage


    % ------------------------------------
    
    \begin{ejercicio}[2 puntos]
        Se realiza un estudio para observar el tiempo que tardan en resolver un problema unos escolares que han seguido un curso de formación por módulos. Se observa el número de módulos que han superado ($X$) junto con el tiempo en minutos que tardan en resolver el problema ($Y$).
        \begin{table}
            
        \end{table}
            \begin{tabular}{|c|ccc|}
                \hline
                $X \setminus Y$ & $\left[1-9\right]$ & $\left]9-21\right]$ & $\left]21-39\right]$ \\ \hline
                $2$ & $0$ & $1$ & $5$ \\
                $4$ & $0$ & $5$ & $5$ \\
                $5$ & $5$ & $3$ & $0$ \\
                $8$ & $15$ & $1$ & $0$ \\ \hline
            \end{tabular}
        \begin{center}
        
        \end{center}

        \begin{enumerate}[label=\alph*)]
            \item Calcular el porcentaje de escolares cuyo tiempo en resolver el problema está comprendido entre la media y la moda de dicha variable.
            \item Calcular la curva de regresión de tipo I de $Y$ en función de $X$.
            \item Estimar el valor de $Y$ cuando $X=4$ mediante una recta de regresión mínimo-cuadrática y dar una medida de la bondad de la predicción. ¿Coincide este valor con la estimación de menor error cuadrático medio de $Y$ cuando $X=4$? Razona la respuesta.
            \item Ajustar a los datos un modelo de regresión hiperbólico para predecir el tiempo de respuesta conociendo el número de módulos superados y predecir el tiempo de respuesta de un estudiante que superó $4$ módulos.
            \item Para predecir el tiempo de respuesta, ¿qué modelo de regresión es más adecuado, el lineal o el hiperbólico?
        \end{enumerate}
    \end{ejercicio}

    \begin{ejercicio}[1 punto]
        Una persona confecciona dos listas de diez artículos para la compra. Al salir de casa elige al azar una de las listas, la primera con probabilidad $0.4$ y la segunda con probabilidad $0.6$. Al llegar al supermercado comprueba que solo lleva dinero para comprar tres artículos y decide elegirlos al azar de los diez de la lista seleccionada.
        \begin{enumerate}[label=\alph*)]
            \item Suponiendo que en el supermercado solo hay existencias de cinco artículos de la primera lista y siete de la segunda, calcular la probabilidad de que la persona encuentre al menos dos de los artículos que ha decidido comprar.
            \item Si encuentra los tres artículos que ha decidido comprar, ¿cuál es la probabilidad de que provengan de la segunda lista?
        \end{enumerate}
    \end{ejercicio}

    \begin{ejercicio}[2.5 puntos]
        Sea $X$ una variable aleatoria con función de densidad
        \[
        f(x) = \begin{cases} 
        a & -1 < x < 0 \\ 
        bx & 0 \leq x < 2 
        \end{cases}
        \]
        \begin{enumerate}[label=\alph*)]
            \item Sabiendo que $E(X)=\nicefrac{5}{12}$, determinar $a$ y $b$.
            \item Dada otra variable aleatoria con media $1$ y varianza $2$, ¿cuál de las dos está más dispersa?
            \item Calcular la función de distribución y la mediana de $X$.
            \item Calcular la función de densidad de $Y$ definida de la siguiente forma:
            \[
            Y = g(X) = \begin{cases} 
            X^2 & -1 < x < 0 \\ 
            X & 0 \leq x < 2 
            \end{cases}
            \]
        \end{enumerate}
    \end{ejercicio}

    \begin{ejercicio}[1.5 puntos]
        En un centro de datos se inspeccionan distintos componentes antes de ponerlos en servicio. Cada jornada se trabaja con uno de los tres lotes siguientes:
        \begin{itemize}
            \item Lote A: $12$ módulos de memoria RAM, de los cuales $5$ son defectuosos.
            \item Lote B: $9$ discos SSD, de los cuales $3$ son defectuosos.
            \item Lote C: $15$ fuentes de alimentación, de las cuales $4$ son defectuosas.
        \end{itemize}
        El responsable selecciona un lote al azar, con probabilidades:
        $$P(A)=0.25; \quad P(B)=0.45; \quad P(C)=0.30.$$
        \begin{enumerate}[label=\alph*)]
            \item Si el responsable elige el lote A y extrae $5$ unidades sin reemplazamiento para una prueba de funcionamiento, ¿cuál es la probabilidad de que exactamente $2$ sean defectuosas?
            \item Si elige el lote B y extrae discos de uno en uno con reemplazamiento hasta obtener el cuarto disco defectuoso, ¿cuál es la probabilidad de que el cuarto defectuoso aparezca en la séptima extracción? Calcular el número medio de discos analizados hasta encontrar el cuarto defectuoso.
            \item El número de errores de conexión detectados por hora en un servidor sigue una distribución de Poisson con media $1.8$. Calcular la probabilidad de que en una hora se detecten exactamente $2$ errores y la de que en tres horas se detecten como mucho $4$ errores.
        \end{enumerate}
    \end{ejercicio}

    \begin{ejercicio}[3 puntos]
        Rodea la opción correcta en cada pregunta.
        
        Puntuación por pregunta: respuesta correcta $+0.5$; incorrecta $-0.2$; en blanco $0$.
        
        \begin{enumerate}[label=\arabic*.]
            \item \textbf{Indica la afirmación correcta:}
            \begin{enumerate}[label=\alph*)]
                \item Si las rectas de regresión ajustadas a una nube son $4x+y=1$; $5x+y=2$, entonces $r_{xy} = \nicefrac{2}{\sqrt{5}}$.
                \item El coeficiente de correlación lineal entre dos variables $X$ e $Y'=-3Y+2$ es el mismo que entre $X$ e $Y$.
                \item Si la varianza de una variable estadística $X$ vale $4$ y el error cuadrático medio asociado a la recta de regresión de $X$ sobre otra variable $Y$ vale $3$, entonces la razón de correlación de $X$ sobre $Y$ es menor que $0.25$.
                \item Si todos los valores de una variable estadística $X$ están comprendidos en un intervalo de amplitud $2$, la varianza de $X$ es menor o igual que $4$.
            \end{enumerate}
            
            \item \textbf{En una variable estadística bidimensional es falso que:}
            \begin{enumerate}[label=\alph*)]
                \item Conociendo las distribuciones marginales están unívocamente determinadas todas las distribuciones condicionadas y la distribución conjunta.
                \item Conociendo la distribución conjunta, están unívocamente determinadas las distribuciones marginales y condicionadas.
                \item Conociendo todas las distribuciones condicionadas para una de las dos variables, están unívocamente determinadas las distribuciones marginales y la distribución conjunta.
                \item Todas las anteriores.
            \end{enumerate}
            
            \item \textbf{Sea $X$ una variable aleatoria con media $m$ distinta de cero y varianza cero. Entonces:}
            \begin{enumerate}[label=\alph*)]
                \item $P(X=0)=1$.
                \item La probabilidad de que $X$ sea distinto de cero es $1$.
                \item $P(X=m)=0$.
                \item Todas las anteriores son falsas.
            \end{enumerate}
            
            \item \textbf{Indica la afirmación correcta:}
            \begin{enumerate}[label=\alph*)]
                \item Si $X \sim B(n,p)$, entonces $P(X = k + 1) = \dfrac{(n-k)p}{(k+1)(1-p)}\cdot P(X = k)$, con $k = 0,1,\dots,n-1$.
                \item La distribución geométrica modela el número de pruebas antes de que ocurra el primer éxito.
                \item Si $X \sim P(\lambda)$, entonces $P(X = k + 1) = \dfrac{\lambda}{k}\cdot P(X = k)$, $k = 0,1,\dots$
                \item Ninguna de las anteriores es cierta.
            \end{enumerate}
            
            \item \textbf{Indica la afirmación correcta:}
            \begin{enumerate}[label=\alph*)]
                \item Dados dos sucesos aleatorios $A$ y $B$, si $P(A)+P(B)>1$, entonces $A$ y $B$ son incompatibles.
                \item Dados dos sucesos aleatorios $A$ y $B$ con $P(B)>0$, si $A$ y $B$ son incompatibles, son independientes.
                \item Si $x$ e $y$ son dos números reales positivos menores o iguales que $a$, la probabilidad de que su suma sea menor o igual que $a$ y su producto sea menor o igual que $b$ es $\nicefrac{1}{2}$ si $a < 2\sqrt{b}$.
                \item Dados dos sucesos aleatorios $A$ y $B$, tales que $P(A)>0$, $P(B)>0$ y $P(A|B)>P(A)$, entonces $P(B|A)=P(B)$.
            \end{enumerate}
            
            \item \textbf{Indica la respuesta correcta:}
            \begin{enumerate}[label=\alph*)]
                \item Los teoremas de cambio de variable solo pueden aplicarse entre variables del mismo tipo.
                \item Toda función continua de $\bb{R}$ en $\bb{R}$ es medible.
                \item Cualquier colección numerable de números reales que sumen uno constituye la función masa de probabilidad de alguna variable aleatoria discreta.
                \item Si una función real de variable real es no negativa e integra uno es la función de densidad de alguna variable aleatoria continua.
            \end{enumerate}
        \end{enumerate}
    \end{ejercicio}

    \newpage
    \setcounter{ejercicio}{0} % Reiniciar contador de ejercicios
    \noindent
    % \textbf{Solución.}

\end{document}